\subsection{Axes IR}
\label{sec:axes-ir}
Axes IR describes an operation's axis connections, output axis order,
and element computation. It provides a common representation for composing
operations, constructing replacement attributes, and checking rewrites.

\subsubsection{Symbolic Axis Groups}
\label{sec:axis-groups}
TEPL binds the symbolic axis groups introduced in
Section~\ref{sec:logical-axes} to operand positions.
While TensorRight's aggregated axes are set-based~\cite{tensorright},
TEPL preserves axis order and operand positions for pairing axes and
constructing attributes. An interface
$\mathsf{I}=\langle a_0,\ldots,a_{n-1}\rangle$ lists a tensor's logical axes in dimension
order. For tensor $t$, we write this interface as $\mathsf{I}_t$.
A zero-based position sequence $q$ locates an ordered group $\vec G$:
\begin{equation}
 \gath(\mathsf{I},q)=\vec G.
 \label{eq:mapping}
\end{equation}
For example, if $\mathsf{I}=\langle i,k,j\rangle$ and $q=[2,0]$,
then $\vec G=\langle j,i\rangle$: positions 2 and 0 select $j$ and $i$,
in that order.


\sidefigure{figures/Symbolic axis groups.pdf}{Symbolic axis groups for generalized dot. Attribute lists bind the batch
($\vec B$), contraction ($\vec K$), and free ($\vec I,\vec J$) groups to operand positions.
Paired axes share coordinates, and summing over $\vec K$ leaves output order
$\vec B\catseq\vec I\catseq\vec J$.}{fig:axis-groups}{Four arbitrary-length groups for batch, left-free, contraction, and right-free
axes are connected to operand position lists and an output-axis construction.}
Each operation defines its Axes IR representation through two blocks:
\code{axes} and \code{semantics}.

The \code{axes} block declares symbolic groups with \code{bind}, which
associates each group with positions in its operands. Binding the same
group in two operands pairs their selected axes in list order.
The block's \code{yield} specifies the output axis order.

The \code{semantics} block describes how each output element is computed
using the groups defined in \code{axes}. For example, a dot multiplies
operand elements at paired coordinates and sums over its contraction group.

Figure~\ref{fig:axis-groups} gives a concrete example of dot's symbolic
axis groups and their operand bindings.

Listing~\ref{lst:dot-declaration} shows how this representation is written
in a TEPL operation declaration. The \code{attrs} block declares \code{lb},
\code{rb}, \code{lc}, and \code{rc} as \code{axis[]} fields storing batch
and contraction positions. The \code{axes} block binds these positions to
logical groups; \code{semantics} computes the output values.


\begin{lstlisting}[language=TEPL,
  caption={Axes IR representation of generalized dot, with explicit attributes, axis bindings, and value semantics.},
  label={lst:dot-declaration}]
op dot(lhs: tensor, rhs: tensor) -> tensor {
  attrs {
    lb, lc, rb, rc: axis[];
  }
  axes {
    let lf = ordered_complement(all_axes(lhs), attrs.lb ++ attrs.lc);
    let rf = ordered_complement(all_axes(rhs), attrs.rb ++ attrs.rc);
    bind B { lhs: attrs.lb, rhs: attrs.rb };
    bind K { lhs: attrs.lc, rhs: attrs.rc };
    bind I { lhs: lf };
    bind J { rhs: rf };
    yield B ++ I ++ J;
  }
  semantics {
    yield Reduce(K, Product(lhs, rhs));
  }
}
\end{lstlisting}
Comma-separated names abbreviate separate attribute declarations.
In Listing~\ref{lst:dot-declaration}, \code{all_axes} lists operand positions,
and \code{ordered_complement(a,b)} removes from \code{a} all entries
appearing in \code{b}, preserving the order of the remaining entries.
For example, removing \code{[0,1]} from \code{[2,0,3,1]} gives \code{[2,3]}.
The bindings pair batch and contraction axes using the attribute
lists. The \code{axes} block yields the output order \code{B ++ I ++ J},
while \code{semantics} multiplies operand elements and sums over \code{K}.
This representation allows a single operation definition to support
different tensor ranks. Appendix~\ref{app:axes} provides a concrete layout
and primitive reference tables.
\subsubsection{Tensor Computation in the \texorpdfstring{\code{semantics}}{semantics} Block}
\label{sec:value-expressions}
The \code{semantics} block represents an operation using basic mathematical
primitives, including operand reads, scalar arithmetic, and summation.
Ordinary operand reads use coordinates established by the \code{axes}
bindings. Summation ranges over a symbolic group while holding output
coordinates fixed. The \code{yield} statement specifies the scalar
expression for each output element. Combining these primitives lets users
define tensor computations through scalar arithmetic and summation over
symbolic axis groups. Table~\ref{extended:tab:value-expressions} in
Appendix~\ref{app:axes} lists the core primitives.

TEPL also supports explicit operand reads through \code{Read(input,indices)},
with index expressions computing access coordinates for operations such
as slices and convolutions. These features extend operation representation,
but synthesis involving general index expressions is outside the scope
of our completeness result. Appendix~\ref{app:axes} documents the index
primitives and gives an access example.
